\documentclass[11pt]{article}
\usepackage[margin=1in]{geometry}
\usepackage{amsmath, amssymb}
\usepackage{graphicx}
\usepackage{hyperref}
\usepackage{booktabs}

\title{A Demo of the \texttt{session-paper} OpenCode Plugin:\\ Automatic Session Distillation and Paper Generation}
\author{Author placeholder} % TODO: replace with real author name/affiliation
\date{\today}

\begin{document}

\maketitle

\begin{abstract}
We present a basic demo and test of \texttt{session-paper}, an OpenCode server plugin
registered at \texttt{\textasciitilde/.config/opencode/plugins/session-paper.ts}. The plugin
injects a standing directive into every primary session's system prompt instructing the
active agent to distill the session and spawn the \texttt{latex-paper-writer} subagent at
the end of every response, naming the output document from a slugified form of the session
title, and it surfaces TUI toasts when the spawn is detected and when the child paper
session finishes. A 14-check harness exercising the plugin's three hooks passed in full,
and the live end-to-end path was then exercised by spawning the subagent from a primary
session, producing this document. We report the design under test, the harness checks, and
the observed end-to-end results.
\end{abstract}

\section{Introduction}
\label{sec:intro}

Large language model agent sessions often produce useful structure---decisions, methods,
and findings---that is lost when the conversation ends. Motivated by the desire to turn
completed sessions into durable, citable documents, the \texttt{session-paper} plugin
automates the last step of a session: distilling it into a LaTeX paper without requiring
the user to prompt for this manually.

The plugin operates through three mechanisms working together: (i)~a standing directive
injected into the system prompt of primary sessions, (ii)~detection of the moment the
active agent spawns the \texttt{latex-paper-writer} subagent, and (iii)~lifecycle tracking
of the resulting child paper session, with user-visible TUI toasts marking the spawn and
the child's completion.

The research question for this demo is straightforward: \emph{does the plugin function as
designed across its three hooks, and does the automatic paper-generation loop complete end
to end in a live session?} We answer it with a scripted test harness plus a live
end-to-end run whose output is this very document.

\section{Background and Related Work}
\label{sec:background}

The plugin builds on OpenCode's extension points: registration in the OpenCode
configuration, hooking of chat system-prompt transformation, interception of tool
execution, and subscription to session events. Related in spirit are prior efforts to
auto-generate scholarly artifacts from agent transcripts, but our scope here is narrower
and strictly empirical: verify the plugin's own contract, not the broader quality of
model-generated papers. No external literature was surveyed for this demo; references are
limited to the artifacts actually exercised.

\section{Methods / Approach}
\label{sec:methods}

\subsection{Plugin under test}
The plugin file \texttt{\textasciitilde/.config/opencode/plugins/session-paper.ts} was
registered in the OpenCode configuration and confirmed active: its directive appeared in
the session's system prompt, with the slug \texttt{testing-session-paper-plugin-demo}
derived from the session title ``Testing session-paper plugin demo.''

The plugin's three hooks are:

\begin{enumerate}
  \item \texttt{experimental.chat.system.transform} --- injects the standing directive
        into primary sessions only.
  \item \texttt{tool.execute.after} --- raises an info toast when a task tool spawns a
        subagent with \texttt{subagent\_type} equal to \texttt{latex-paper-writer}.
  \item An event handler --- tracks \texttt{session.created}/\texttt{session.updated}
        titles, links the child paper session to its parent, and raises a success toast on
        \texttt{session.idle}.
\end{enumerate}

\subsection{Test harness}
A 14-check \texttt{bun test} harness with a mock client exercised all three hooks. Beyond
the happy path, the harness asserted negative and edge behaviors:

\begin{itemize}
  \item the directive is suppressed in child/subagent sessions, preventing recursion;
  \item unrelated tools spawn no toast;
  \item unknown sessions fall back to \texttt{client.session.get};
  \item renamed sessions update the derived slug;
  \item idle toasts fire exactly once per child session.
\end{itemize}

\subsection{Live end-to-end run}
Finally, the real path was exercised: the \texttt{latex-paper-writer} subagent was spawned
from the primary session, the plugin detected the spawn and toasted, and the child session
produced this document.

\section{Results / Findings}
\label{sec:results}

All 14 harness checks passed, covering directive injection, slug derivation,
child-session gating, and the toast lifecycle. Table~\ref{tab:hooks} summarizes the
hook-level outcomes.

\begin{table}[h]
  \centering
  \begin{tabular}{@{}llc@{}}
    \toprule
    Hook & Behavior verified & Outcome \\
    \midrule
    \texttt{system.transform} & Directive injected into primary sessions only & Pass \\
    \texttt{tool.execute.after} & Toast on \texttt{latex-paper-writer} spawn only & Pass \\
    Event handler & Title tracking, child linking, single idle toast & Pass \\
    \bottomrule
  \end{tabular}
  \caption{Hook-level results from the 14-check harness.}
  \label{tab:hooks}
\end{table}

The live end-to-end path also succeeded: the subagent spawn was detected and toasted, and
this document was written to disk as \texttt{testing-session-paper-plugin-demo.pdf} with
source \texttt{testing-session-paper-plugin-demo.tex}.

\section{Discussion}
\label{sec:discussion}

The results indicate the plugin works as designed. The negative checks matter as much as
the positive ones: suppressing the directive in child sessions prevents a paper-writing
subagent from recursively spawning further paper writers, and gating toasts on the
specific \texttt{latex-paper-writer} subagent type avoids noise from unrelated tool calls.
The rename check confirms slug derivation tracks live title edits rather than a stale
snapshot, and the exactly-once idle toast prevents duplicated completion notifications.

This demo deliberately limits its claims. It verifies the plugin's contract and one live
run; it does not evaluate paper quality across sessions, robustness under concurrent
sessions, or behavior when the subagent fails. Those remain open for future testing.

\section{Conclusion}
\label{sec:conclusion}

The \texttt{session-paper} plugin is operational: directive injection, slug derivation,
child-session gating, and the toast lifecycle were all verified by a 14-check harness, and
the automatic paper-generation loop functions end to end, as demonstrated by the
production of this document.

\section*{Acknowledgments}
% TODO: optional acknowledgments.

\begin{thebibliography}{9}
\bibitem{session-paper-plugin}
\texttt{session-paper} OpenCode server plugin source,
\texttt{\textasciitilde/.config/opencode/plugins/session-paper.ts}; exercised in this
session via its configuration registration and 14-check \texttt{bun test} harness.

\bibitem{session-paper-output}
This document itself: \texttt{testing-session-paper-plugin-demo.tex} /
\texttt{testing-session-paper-plugin-demo.pdf}, generated by the
\texttt{latex-paper-writer} subagent spawned through the plugin's end-to-end path.
\end{thebibliography}

\end{document}
